\documentclass[a4paper,11pt]{article}
\usepackage[T1]{fontenc}
\usepackage[utf8]{inputenc}
\usepackage{booktabs}
\usepackage{graphicx}
\usepackage[hidelinks]{hyperref}

% Standalone draft of Sections 3.2-3.5 (additions), 4 and 5 for the DEXA paper. In the paper,
% drop the preamble and the bibliography and paste the sections; table labels stay valid.
% Numbers: results/dgx on branch QED_LEAF of vaibhav-baj/llm-smarthome-extension (runs of
% 23-25 October 2026). Rows marked "pending" are still running.

\begin{document}

\title{Draft: system description, benchmark and evaluation}
\date{}
\maketitle

% =============================================================================
\section{LLM-based multi-agent system (Sections~3.2--3.5)}
% =============================================================================

\subsection{Current system}
\label{sec:current-system}

The current system keeps the two-level design of the original architecture, centralized planning
and device-level execution, and adds five mechanisms around it. (1)~Every device agent publishes
an \emph{agent card}: for each operation it offers, the card declares the parameters with their
types, ranges and defaults, the state variables the operation reads and writes, its preconditions
and its postconditions. The cards are the single source of truth for planning, validation, the
allocation checker and the external benchmark adapters. (2)~Tasks carry explicit dependencies, so
independent tasks run in parallel and dependent ones wait for verified results. (3)~Every planned
action is validated against its card before dispatch and its outcome is checked against the
observed home state afterwards. (4)~Failures are repaired locally where possible and replanned
otherwise. (5)~A batch of commands executes only if the formally verified allocation checker of
Section~\ref{sec:lean} admits it. All five can be switched off individually, which gives the
ablations of Section~\ref{sec:evaluation}. The execution loop is: intent analysis (optional),
task planning, domain planning by the device agents, plan validation, admission and scheduling of
one round of independent commands, execution with outcome checks, and completion of the round,
repeated until every task is finished or terminally failed.

\subsection{Task planning}
\label{sec:task-planning}

The task planner receives the user request, the agent cards reduced to their operations and
preconditions, the rooms, the current public state and, when enabled, the intent analysis. It
returns a task queue in which every task names its agent, a plain-language goal, the devices or
rooms it may touch (\emph{scope}), its constraints (times, quantities, conditions, negation) and
the tasks it depends on. A task may already contain concrete actions, otherwise the device agent
refines it. Sub-requests that cannot be fulfilled with the available devices are listed as
rejected rather than planned. The planner is instructed to plan the entire request as one task
per distinct goal and to use only agent and device identifiers from the cards; a plan that names
an unknown agent, duplicates a task identifier or contains a dependency cycle is rejected and the
planner is asked once more with the error as context. The device agents then turn each task into
concrete operations from their own card, declaring which earlier results an argument is taken
from (\emph{argument sources}); \texttt{validate\_plan} resolves these references, inserts
explicitly approved prerequisite operations (for example switching a device on when a
precondition requires it) and rejects plans with contradictory goals on the same device.

\subsection{Dynamic collaboration mechanism}
\label{sec:collaboration}

Collaboration is no longer a free-form message between agents. An agent that needs information
owned by another device declares an \emph{information task} with the owning agent as its
parent, or references a predecessor's output field as an argument source. Both become edges of
the dependency graph, so the scheduler orders them, the verified result of the predecessor is
passed on with its observation time and version, and the allocation checker refuses to run a
dependent command whose evidence has expired or whose source state has changed since it was
observed (Section~\ref{sec:lean}, condition~3). Results of completed tasks remain available to
later tasks through the shared task history, which preserves the implicit information reuse of
the original system. When a device agent cannot complete a task, repair is local first: only the
actions that depend on the failed one, by task dependency or by reading a state variable it
writes, are skipped and replanned by that agent (the soundness of this rule is proved in
\texttt{Repair.lean}); a failure that crosses agents is handed back to the task planner, which
replans the affected tasks only, preserving their identifiers and goals.

\subsection{Formally verified allocation checking: auditing the proofs}
\label{sec:lean}

% (The existing Section 3.5 text stays as it is; this paragraph is appended after
%  "... implemented in Python and not verified.")
That the proofs compile is not by itself evidence that the stated theorems hold: the Lean
compiler accepts any valid proof of any theorem, and it does not know which theorem was intended,
which axioms a project permits, or whether a definition behind an unchanged statement has been
weakened. We therefore audit the proofs with QED\_LEAF, a fail-closed checker for Lean
developments that re-elaborates the targets and compares them with a protected copy of the
specification. For each target theorem it checks that the proof contains no unproven step
(\texttt{sorry}), that the statement is structurally identical to the protected reference
(statement match), that every definition the statement depends on is restated identically
(definition integrity), that only the permitted axioms \texttt{propext} and \texttt{Quot.sound}
are used (axiom policy), and that the theorem does not hold vacuously. The audit runs before
every evaluation that uses the checker and its outcome is recorded with the run. It takes about
3.5 minutes on a CPU, is cached while the Lean sources, policy and reference are unchanged, and
is not part of request handling, so task scores and latency are unaffected.

Table~\ref{tab:qedleaf-evidence} contrasts the audit with compilation alone on four versions of
the proofs. All four compile. The audit passes only the current version; it rejects an earlier
version whose soundness and completeness proofs depended on \texttt{Classical.choice} through a
single \texttt{simp} step, a policy violation that compilation had not surfaced
(Table~\ref{tab:qedleaf-axioms}), and it rejects two deliberately weakened versions, one with a
weaker theorem published under the protected name and one with an independence definition that
drops the exclusive-resource condition. Like the checker's guarantees, the audit covers only the
Lean development; the Python translation remains outside its scope.

\begin{table}[htbp]
  \centering
  \small
  \caption{Compilation alone vs.\ the QED\_LEAF audit on four versions of the allocation-checker
    proofs (both with Lean~4.33.1; QED\_LEAF analyzer revision 9262eb0). Every version compiles;
    only the audit rejects the three flawed ones. Audit targets:
    \texttt{checkAllocation\_sound}, \texttt{checkAllocation\_complete} and
    \texttt{GraphWF.acyclic}, compared against a protected reference specification; permitted
    axioms \texttt{propext} and \texttt{Quot.sound}.}
  \label{tab:qedleaf-evidence}
  \resizebox{\linewidth}{!}{%
  \begin{tabular}{llcll}
    \toprule
    Proof version & Change & Compiles & QED\_LEAF audit & Failed checks \\
    \midrule
    Upstream, before the fix & one \texttt{simp} step reaches \texttt{Classical.choice} & yes & \textbf{failed} & axiom policy (soundness, completeness) \\
    Current, after the fix   & \texttt{beq\_iff\_eq.mpr} instead of \texttt{simp}        & yes & \textbf{passed} (24/24 checks) & none \\
    Negative control 1       & weaker theorem under the name \texttt{checkAllocation\_sound} & yes & \textbf{failed} & statement match (soundness) \\
    Negative control 2       & independence definition without the exclusive-resource condition & yes & \textbf{failed} & definition integrity (soundness, acyclicity) \\
    \bottomrule
  \end{tabular}}
\end{table}

\begin{table}[htbp]
  \centering
  \small
  \caption{Axioms of the soundness and completeness theorems as reported by Lean
    (\texttt{\#print axioms}). Both versions compile; the audit's axiom policy permits only
    \texttt{propext} and \texttt{Quot.sound}.}
  \label{tab:qedleaf-axioms}
  \resizebox{\linewidth}{!}{%
  \begin{tabular}{llc}
    \toprule
    Proof version & Axioms & Axiom policy \\
    \midrule
    Upstream, before the fix & \texttt{Classical.choice}, \texttt{propext}, \texttt{Quot.sound} & failed \\
    Current, after the fix   & \texttt{propext}, \texttt{Quot.sound}                           & passed \\
    \bottomrule
  \end{tabular}}
\end{table}

% =============================================================================
\section{Benchmark}
\label{sec:benchmark}
% =============================================================================

The revised edition of our benchmark contains 127 scenarios: a \emph{core} suite of 104 and a
\emph{robustness} suite of 23. The core suite keeps 57 of the 60 original cases, revised where
their annotations were ambiguous (three cases with subjective scene assumptions were dropped and
one replaced), and adds 40 cases designed along four dimensions that the original set hardly
covers: dependency structure (chains, fan-out, fan-in and conditional branches of depth up to
three; 15 cases), ambiguity (vague requests with several acceptable interpretations; 9),
scale (five to eight intents across up to seven device types; 8) and concurrency
(contradictory or parallelizable requests for the same devices; 8). Fourteen further cases are
\emph{variants} of existing ones that change the home state (a busy calendar, an empty fridge,
rain) so that the correct action differs although the request is the same. Each case carries a
complexity vector (number of intents and device types, dependency depth and width, ambiguity,
temporal and spatial constraints, conflicts) and an \emph{expected behavior}: 101 cases must be
executed, and 26 require a focused clarification question about an unresolved part of the
request while the clear, independent goals are still carried out. The robustness suite injects
device failures into 16 scenarios, each with an unperturbed control: a device is unavailable
(11), answers only after a transient failure (5), reports success without effect (1), or its
state is changed externally after planning (1). A correct response completes the independent
goals, reports the failure truthfully, replans dependent goals instead of executing them with
missing information, and, for transient failures, recovers the requested goal.

Every case runs in a fresh simulated home with versioned initial values, a virtual clock and
synthetic catalogues for media, recipes and search results, and is scored automatically: there
is no manual Complete/Partial/Failed step and no LLM judge. The evaluator checks the required
end-state predicates, that information goals rest on an actual recorded observation, that result
dependencies were satisfied before dependent calls, that nothing outside the allowed changes was
written, that existing calendar entries and list items were preserved, that scheduled effects
ran at the right virtual time with the right arguments, and, for clarification cases, that a
question was asked and the unresolved parts were left untouched. Free-text answers are checked
lexically against published term lists, which is deliberately weak but judge-free; the binary
score therefore means ``passed the published deterministic checks''. The benchmark exposes one
episode interface (request, public context, tool schemas, an \texttt{invoke} function) that every
compared system uses, so no system sees gold goals, failure labels or annotations. The original
60-case edition remains available unchanged for comparison with earlier results.

% =============================================================================
\section{Evaluation}
\label{sec:evaluation}
% =============================================================================

All experiments use Gemma~4 31B~\cite{gemma4} served locally through Ollama (4-bit quantization,
Q4\_K\_M) with thinking disabled, temperature 0 and seed 42, on a DGX-2 with Tesla V100 (32\,GB)
GPUs. Two systems are compared with the same model: our multi-agent system (MAS) with every
mechanism of Section~\ref{sec:current-system} enabled and the Lean checker admitting every
command batch (the QED\_LEAF audit passed before every such run), and a single ReAct
agent~\cite{react} that receives the same episode interface and tools. We evaluate on our
extended benchmark (Section~\ref{sec:eval-own}), the 95 cases of the Home Assistant Assist
dataset (Section~\ref{sec:eval-ha}), the 600 episodes of SimuHome (Section~\ref{sec:eval-simuhome})
and a stratified sample of 1,000 HomeBench test instructions drawn proportionally to instruction
type with seed 42 (Section~\ref{sec:eval-homebench}). Each benchmark is scored by its native
evaluator. Table~\ref{tab:systems} summarizes the results.

\begin{table}[htbp]
  \centering
  \small
  \caption{Both systems with the same Gemma~4 31B on four benchmarks. Success rates are each
    benchmark's native binary score; HomeBench reports exact match and operation-level F1.
    The MAS runs with the Lean checker; ReAct has no admission step. ``pending'' = run in progress.}
  \label{tab:systems}
  \resizebox{\linewidth}{!}{%
  \begin{tabular}{llrrrrr}
    \toprule
    Benchmark & Metric & Cases & MAS + Lean & ReAct & Difference (MAS $-$ ReAct) & 95\% CI \\
    \midrule
    DEXA extended (core + robustness) & success rate & 127 & 0.291 & \textbf{0.543} & $-0.252$ & $[-0.354, -0.150]$ \\
    Home Assistant Assist             & success rate & 95  & 0.842 & 0.863 & $-0.021$ & n.s.\ (2 cases) \\
    SimuHome                          & success rate & 600 & 0.417 & pending & -- & -- \\
    HomeBench                         & exact match / F1 & 1,000 & 0.649 / 0.703 & pending & -- & -- \\
    \bottomrule
  \end{tabular}}
\end{table}

% -----------------------------------------------------------------------------
\subsection{Own benchmark}
\label{sec:eval-own}
% -----------------------------------------------------------------------------

On the extended benchmark the MAS passes 37 of 127 cases (0.291) and the ReAct agent 69 (0.543);
on the 127 shared cases the difference of $-0.252$ has a 95\% bootstrap confidence interval of
$[-0.354, -0.150]$, so with this model the single agent is better. Table~\ref{tab:dexa-breakdown}
locates the difference. It lies entirely in the core suite and grows with complexity: on simple
cases the gap is 0.20, on moderate and complex cases 0.30. By dimension it is largest where the
MAS's explicit structure should help most, on dependency chains (0.13 against 0.67), and absent on
ambiguity. On the robustness suite the two systems tie at 0.348, and the MAS is the only one that
recovers from transient device failures (0.40 against 0.00) and the one \emph{no-effect} case,
which is what its outcome checks and local repair are designed for. Both systems handle the 26
clarification cases poorly (0.15 and 0.19). The MAS ended 17 cases with a system error, the ReAct
agent none. The allocation checker rejected no batch in any of the 127 cases: every batch the
planner proposed was valid, so the checker neither cost nor gained a case here.

\begin{table}[htbp]
  \centering
  \small
  \caption{Extended benchmark: success rate by suite, category, dimension and expected behavior
    (same model; automatic evaluator). The original dimension holds the 57 revised original cases
    and the variants.}
  \label{tab:dexa-breakdown}
  \begin{tabular}{llrrr}
    \toprule
    Group & & $n$ & MAS + Lean & ReAct \\
    \midrule
    Suite       & core                 & 104 & 0.279 & 0.587 \\
                & robustness           & 23  & 0.348 & 0.348 \\
    \midrule
    Category    & simple               & 61  & 0.410 & 0.607 \\
                & moderate             & 37  & 0.189 & 0.486 \\
                & complex              & 29  & 0.172 & 0.483 \\
    \midrule
    Dimension   & original             & 64  & 0.250 & 0.578 \\
                & dependency structure & 15  & 0.133 & 0.667 \\
                & ambiguity            & 9   & 0.444 & 0.444 \\
                & scale                & 8   & 0.375 & 0.625 \\
                & concurrency          & 8   & 0.500 & 0.625 \\
                & failure recovery     & 23  & 0.348 & 0.348 \\
    \midrule
    Behavior    & execute              & 101 & 0.327 & 0.634 \\
                & clarify              & 26  & 0.154 & 0.192 \\
    \midrule
    Failure     & none (incl.\ controls) & 109 & 0.275 & 0.587 \\
                & unavailable          & 11  & 0.273 & 0.364 \\
                & transient            & 5   & 0.400 & 0.000 \\
    \bottomrule
  \end{tabular}
\end{table}

Where the loss arises is visible from a classification probe (Table~\ref{tab:classification}).
Asked directly which device types a request involves, Gemma~4 answers the 115 extended cases
with device labels almost perfectly (macro-F1 0.957, recall 0.992). On the 60 original cases,
where the end-to-end MAS can be scored the same way, the agents it actually assigned a task reach
a macro-F1 of only 0.638 with recall 0.497: the devices it activates are almost always appropriate
(precision 0.948), but it activates too few of them. The model understands the request; the task
planner drops goals when there are many, and a dropped goal fails the case under the binary
score. The ReAct agent, which never decomposes the request but works through it one tool call at
a time, does not have this failure mode.

% -----------------------------------------------------------------------------
\subsection{Benchmark 1: Home Assistant dataset}
\label{sec:eval-ha}
% -----------------------------------------------------------------------------

Both systems fulfil most of the 95 requests: the MAS 80 (0.842), the ReAct agent 82 (0.863), a
difference of two cases. The MAS needs 107\,s per case against 17\,s, since every request passes
through planning, validation and outcome checks. Table~\ref{tab:gemma4-homeassistant} breaks the
MAS result down by device domain: covers, fans, to-do lists and valves succeed in every case,
whereas lights, locks and vacuums, whose requests more often address several entities or areas at
once, lie between 0.67 and 0.71.

\begin{table}[htbp]
  \centering
  \small
  \caption{MAS + Lean on the Home Assistant Assist dataset by device domain.}
  \label{tab:gemma4-homeassistant}
  \begin{tabular}{lrr}
    \toprule
    Domain        & $n$ & Success rate \\
    \midrule
    Cover         & 11  & 1.00 \\
    Fan           & 6   & 1.00 \\
    Light         & 13  & 0.69 \\
    Lock          & 9   & 0.67 \\
    Media player  & 21  & 0.86 \\
    To-do list    & 9   & 1.00 \\
    Vacuum        & 17  & 0.71 \\
    Valve         & 9   & 1.00 \\
    \midrule
    All           & 95  & 0.842 \\
    \bottomrule
  \end{tabular}
\end{table}

% -----------------------------------------------------------------------------
\subsection{Benchmark 2: SimuHome}
\label{sec:eval-simuhome}
% -----------------------------------------------------------------------------

Table~\ref{tab:gemma4-simuhome} reports the MAS success rate by query type. State questions
(QT1) are answered correctly in about four of five episodes, whether feasible or not. Workflow
scheduling (QT4) is the weakest type, as the benchmark's authors also observe for other
models~\cite{simuhome}: feasible workflows succeed in 0.26--0.44 of the episodes, and the system
almost never recognizes an infeasible delayed or conditional workflow (QT4-2 and QT4-3 infeasible:
one success each in 50 episodes), answering as if the request could be scheduled.

\begin{table}[htbp]
  \centering
  \small
  \caption{MAS + Lean success rate on SimuHome by query type (50 episodes per cell). QT1 and all
    infeasible episodes are judged by a three-vote LLM panel (Gemma~4 31B, 49,152-token context;
    1,008 judge calls, no errors or truncated prompts).}
  \label{tab:gemma4-simuhome}
  \begin{tabular}{lrrr}
    \toprule
    Query type & Feasible & Infeasible & All   \\
    \midrule
    QT1        & 0.82     & 0.80       & 0.81  \\
    QT2        & 0.42     & 0.34       & 0.38  \\
    QT3        & 0.42     & 0.80       & 0.61  \\
    QT4-1      & 0.26     & 0.26       & 0.26  \\
    QT4-2      & 0.44     & 0.02       & 0.23  \\
    QT4-3      & 0.40     & 0.02       & 0.21  \\
    \midrule
    All        & 0.460    & 0.373      & 0.417 \\
    \bottomrule
  \end{tabular}
\end{table}

% -----------------------------------------------------------------------------
\subsection{Benchmark 3: HomeBench}
\label{sec:eval-homebench}
% -----------------------------------------------------------------------------

On the 1,000-case sample the MAS reaches an exact match of 0.649 and an operation-level F1 of
0.703 (Table~\ref{tab:gemma4-homebench}). Single valid instructions and instructions for
non-existent devices are handled well (exact match 0.848 and 0.905), instructions for unsupported
attributes less so (0.696). Multi-step instructions are the weakest type: most of their individual
operations are correct (F1 0.645), but only 0.098 of them are fully correct, since one wrong or
omitted operation fails the whole instruction; of the 1,509 expected operations in these cases,
32\% were missing and 39\% of the predicted ones were wrong or superfluous. The ReAct run on the
same sample is in progress.

\begin{table}[htbp]
  \centering
  \small
  \caption{MAS + Lean on HomeBench by instruction type (stratified sample of 1,000 test cases, seed 42).}
  \label{tab:gemma4-homebench}
  \begin{tabular}{lrrrrr}
    \toprule
    Instruction type   & $n$   & Exact match & Precision & Recall & F1    \\
    \midrule
    Normal             & 356   & 0.848 & 0.791 & 0.916 & 0.849 \\
    Missing device     & 242   & 0.905 & 0.818 & 0.909 & 0.861 \\
    Missing attribute  & 148   & 0.696 & 0.656 & 0.696 & 0.675 \\
    Multi-step         & 254   & 0.098 & 0.615 & 0.679 & 0.645 \\
    \midrule
    All                & 1,000 & 0.649 & 0.668 & 0.742 & 0.703 \\
    \bottomrule
  \end{tabular}
\end{table}

% -----------------------------------------------------------------------------
\subsection{Performance by intent complexity}
\label{sec:eval-complexity}
% -----------------------------------------------------------------------------

Across benchmarks the same pattern recurs: both systems handle single, explicit requests well, and
performance falls with the number of goals a request contains. On our benchmark the MAS drops from
0.41 (simple) to 0.17 (complex) and the ReAct agent from 0.61 to 0.48; on HomeBench the MAS
falls from 0.85 exact match on single instructions to 0.10 on multi-step ones while still getting
two thirds of the individual operations right; on SimuHome, workflow scheduling is the weakest
query type for both the MAS and the models reported by the benchmark's authors. The classification
probe shows that this is not a limit of the model's understanding, which identifies the involved
devices almost perfectly, but of turning that understanding into a complete plan. For the MAS the
loss is concentrated in task decomposition, and it is larger than for the single agent because a
goal the planner omits is never recovered, whereas the single agent revisits the request after
every tool call. The one place where the MAS's additional structure pays off is failure recovery:
under transient device failures its outcome checks and local repair recover the requested goals
where the single agent does not.

% -----------------------------------------------------------------------------
\subsection{Limitations of the evaluation}
% -----------------------------------------------------------------------------

Several aspects of the setup qualify the results. (1) Each configuration was run once, with a
quantized model and deterministic decoding; we report no variance across seeds. (2) The ReAct
baseline runs without the allocation checker, so the comparison measures task success, not
safety: the checker's guarantees (Section~\ref{sec:lean}) hold only for the MAS. On the extended
benchmark the checker rejected no batch, so it did not affect the MAS score either. (3) SimuHome
scores QT1 and all infeasible episodes with a three-vote LLM judge. We used Gemma~4 31B as the
judge, so the model under test judges its own answers, which may favor it; the judge received
complete prompts (49,152-token context, no truncation). The benchmark's default judge
configuration (Gemma~2, 8,192 tokens) cannot read the judge prompts, which exceed 28,000 tokens,
and scored every answer as a failure in a first run that we discarded. (4) The Home Assistant
results are from a rerun after correcting a harness defect: the adapter's routing argument
\texttt{service\_id}, which is not part of Home Assistant's own tool schema, was required but not
defaulted, so 81 of 95 MAS cases failed before planning could complete. (5) HomeBench was
evaluated on a 1,000-case sample, not the full test set of 17,366 instructions. (6) Our
benchmark's free-text checks are lexical; an answer phrased outside the published term lists
fails although it may be correct, and this affects both systems alike. (7) The checker's
guarantees and the QED\_LEAF audit cover the Lean development; the translation of agent cards,
tasks and observed state into the checker's input is unverified Python.

\begin{table}[htbp]
  \centering
  \small
  \caption{Classification probe: which device types does a request involve? Gemma~4 zero-shot
    against a fine-tuned or zero-shot ModernBERT-large encoder on one task per benchmark
    (macro-averaged; 95\% bootstrap CI of macro-F1). The ``end to end'' row scores the device
    types of the agents the MAS actually assigned a task.}
  \label{tab:classification}
  \resizebox{\linewidth}{!}{%
  \begin{tabular}{llrrrrrc}
    \toprule
    Task & Model & $n$ & Acc. & Precision & Recall & F1 & 95\% CI (F1) \\
    \midrule
    T1 DEXA device types, extended edition & Gemma~4, zero-shot & 115 & 0.817 & 0.926 & 0.992 & \textbf{0.957} & [0.936, 0.975] \\
    T1 DEXA device types, extended edition & ModernBERT, zero-shot & 115 & 0.330 & 0.852 & 0.720 & 0.759 & [0.719, 0.794] \\
    T1 DEXA device types, original edition & Gemma~4, zero-shot & 60 & 0.550 & 0.940 & 0.834 & \textbf{0.872} & [0.825, 0.906] \\
    T1 DEXA device types, original edition & Gemma~4 + MAS, end to end & 60 & 0.267 & 0.948 & 0.497 & 0.638 & [0.552, 0.706] \\
    T1 DEXA device types, original edition & ModernBERT, zero-shot & 60 & 0.267 & 0.909 & 0.510 & 0.611 & [0.524, 0.681] \\
    T2 HomeBench validity (3 classes) & Gemma~4, zero-shot & 1,500 & 0.873 & 0.890 & 0.873 & 0.872 & [0.856, 0.888] \\
    T2 HomeBench validity (3 classes) & ModernBERT, fine-tuned & 1,500 & 0.972 & 0.974 & 0.972 & \textbf{0.972} & -- \\
    T3 SimuHome feasibility (2 classes) & Gemma~4, zero-shot & 600 & 0.592 & 0.607 & 0.592 & 0.577 & [0.537, 0.616] \\
    T3 SimuHome feasibility (2 classes) & ModernBERT, 5-fold CV & 600 & 0.788 & 0.794 & 0.788 & \textbf{0.787} & [0.755, 0.820] \\
    T4 Home Assistant domain (8 classes) & Gemma~4, zero-shot & 95 & 1.000 & 1.000 & 1.000 & \textbf{1.000} & [1.000, 1.000] \\
    T4 Home Assistant domain (8 classes) & ModernBERT, zero-shot & 95 & 0.937 & 0.972 & 0.956 & 0.958 & [0.922, 0.986] \\
    \bottomrule
  \end{tabular}}
\end{table}

% =============================================================================
\section{Conclusion}
% =============================================================================

We extended an LLM-based multi-agent smart home system with device contracts, a dependency graph
with parallel execution, validation with local repair, and an allocation checker whose soundness
and completeness are proved in Lean~4 and whose proofs are audited by QED\_LEAF against a protected
specification. We also extended the benchmark to 127 automatically scored scenarios with a
robustness suite, and evaluated the system with an open 31B model on four benchmarks against a
single ReAct agent with the same model. The verified checker works as intended: it admitted every
valid batch, the audit caught a real axiom-policy violation that compilation had not, and the
proofs now rely only on Lean's standard axioms. The effectiveness result is less favorable: on our
benchmark the single agent passes 0.54 of the cases and the multi-agent system 0.29, and the gap
grows with the number of goals in a request; on Home Assistant the two are within two cases. The
multi-agent system is better only under device failures, where its outcome checks and local repair
recover goals that the single agent loses. The cause is task decomposition, not model
understanding: the model identifies the devices a request involves almost perfectly, and the
planner drops goals when there are many. The architecture's safety argument therefore stands,
while its effectiveness argument depends on fixing decomposition.

% =============================================================================
\section{Future work}
% =============================================================================

The most direct next step is a plan-completeness check: after planning, compare the goals the
model itself extracts from the request with the tasks in the queue, and send the uncovered goals
back to the planner before execution. We also plan to run the multi-agent system without the
allocation checker on the extended benchmark to report the checker's runtime cost precisely, to
repeat all runs over several seeds, to replace the self-judging SimuHome configuration with an
independent judge model, and to complete the ReAct baseline on SimuHome and HomeBench. For the
benchmark, no case yet has a Lean rejection as its correct outcome; a safety suite of conflicting,
out-of-range and stale-evidence requests would evaluate the verified layer itself rather than the
language model. Finally, the fine-tuned encoder's advantage on validity and feasibility
classification (Table~\ref{tab:classification}) suggests a fast pre-check in front of the planner
for requests that name a missing device or cannot be fulfilled, the cases where both systems are
weakest.

% =============================================================================
% Bibliography (standalone only; the paper has its own)
% =============================================================================
\begin{thebibliography}{9}
\bibitem{gemma4} Gemma Team, Google DeepMind: Gemma 4 technical report (2026),
  \url{https://arxiv.org/abs/2607.02770}
\bibitem{react} Yao, S., Zhao, J., Yu, D., Du, N., Shafran, I., Narasimhan, K., Cao, Y.: ReAct:
  Synergizing reasoning and acting in language models. In: The Eleventh International Conference
  on Learning Representations (2023). \url{https://arxiv.org/abs/2210.03629}
\bibitem{homebench} Li, S., Guo, Y., Yao, J., Liu, Z., Wang, H.: HomeBench: Evaluating LLMs in
  smart homes with valid and invalid instructions across single and multiple devices. In:
  Proceedings of the 63rd Annual Meeting of the Association for Computational Linguistics
  (Volume 1: Long Papers), pp. 12230--12250 (2025). \url{https://aclanthology.org/2025.acl-long.597/}
\bibitem{simuhome} Seo, G., Yang, J., Pyo, J., Kim, N., Lee, J., Jo, Y.: SimuHome: A temporal- and
  environment-aware benchmark for smart home LLM agents. In: The Fourteenth International
  Conference on Learning Representations (2026). \url{https://openreview.net/forum?id=LCS1WsGvha}
\end{thebibliography}

\end{document}
